\documentclass[tikz, border=15pt]{standalone}
\usepackage[utf8]{inputenc}
\usepackage[T1]{fontenc}
\usepackage{roboto-mono}
\renewcommand*\familydefault{\ttdefault}

\usepackage{tikz}
\usetikzlibrary{arrows.meta, positioning, calc, decorations.pathmorphing, shapes.gates.logic.US}

\definecolor{stagecolor}{HTML}{1F6FEB}    % Accent Blue
\definecolor{hazardcolor}{HTML}{D29922}   % Amber / Orange
\definecolor{regcolor}{HTML}{7D8590}      % Muted Slate / Gray
\definecolor{boundarycolor}{HTML}{8C959F} % Dashed Line Gray
\definecolor{arrowgreen}{HTML}{2EA043}    % High-contrast green (GitHub green)
\definecolor{arrowpurple}{HTML}{8250DF}   % Vibrant Purple (GitHub Purple)

% Load diagram building utilities (e.g. \drawblock)
\IfFileExists{diagrams/util.tex}{\input{diagrams/util.tex}}{% TikZ Utilities for Hardware Architecture Diagrams
% Provides reusable styles and macros for drawing modules with stacked IO ports.

\usetikzlibrary{calc}

% Pin label style for ports
\tikzset{
    pinlabel/.style={font=\footnotesize\ttfamily\bfseries, text=stagecolor}
}

% Counters for port row stacking (declared once globally)
\newcount\nL
\newcount\nR
\newcount\nRows
\newcount\iL
\newcount\iR

% Macro: \drawblock
% Draws a rectangular hardware module with title and auto-stacked IO port columns.
% Input ports are stacked vertically on the left; output ports on the right.
% #1: Module identifier prefix (e.g. csr_decode)
% #2: Center X
% #3: Center Y
% #4: Width
% #5: Title (typeset with \texttt)
% #6: Left ports list  {port_name/color, ...}
% #7: Right ports list {port_name/color, ...}
% #8: Bottom ports list {port_name/x_offset/color, ...}
% #9: Custom height override (if empty, auto-computed from row count)
\newcommand{\drawblock}[9]{%
  % Count left and right ports to determine row count
  \nL=0
  \foreach \p in {#6} { \global\advance\nL by 1 }
  \nR=0
  \foreach \p in {#7} { \global\advance\nR by 1 }
  \ifnum\nL>\nR \nRows=\nL \else \nRows=\nR \fi
  \ifnum\nRows<1 \nRows=1 \fi
  % Geometry constants
  \pgfmathsetmacro{\rowH}{0.80}%
  \pgfmathsetmacro{\topH}{0.80}%
  \pgfmathsetmacro{\botH}{0.50}%
  \pgfmathsetmacro{\calcH}{\topH + \nRows * \rowH + \botH}%
  \def\overrideH{#9}%
  \ifx\overrideH\empty
    \pgfmathsetmacro{\H}{\calcH}%
  \else
    \pgfmathsetmacro{\H}{#9}%
  \fi
  \pgfmathsetmacro{\xL}{#2 - (#4)/2}%
  \pgfmathsetmacro{\xR}{#2 + (#4)/2}%
  \pgfmathsetmacro{\yT}{#3 + \H/2}%
  \pgfmathsetmacro{\yB}{#3 - \H/2}%
  % Module box
  \draw[draw=stagecolor, line width=1.5pt, fill=stagecolor, fill opacity=0.15]
    (\xL, \yT) rectangle (\xR, \yB);
  % Title
  \node[anchor=north, font=\bfseries\normalsize, text=stagecolor]
    at (#2, \yT - 0.15) {\texttt{#5}};
  % Left side ports (stacked column, top-to-bottom)
  \iL=0
  \foreach \pname/\pcol in {#6} {
    \global\advance\iL by 1
    \pgfmathsetmacro{\py}{\yT - \topH - (\the\iL - 0.5)*\rowH}%
    \coordinate (#1-\pname) at (\xL, \py);
    \node[pinlabel, anchor=west, text=\pcol] at (\xL + 0.12, \py) {\detokenize\expandafter{\pname}};
  }
  % Right side ports (stacked column, top-to-bottom)
  \iR=0
  \foreach \pname/\pcol in {#7} {
    \global\advance\iR by 1
    \pgfmathsetmacro{\py}{\yT - \topH - (\the\iR - 0.5)*\rowH}%
    \coordinate (#1-\pname) at (\xR, \py);
    \node[pinlabel, anchor=east, text=\pcol] at (\xR - 0.12, \py) {\detokenize\expandafter{\pname}};
  }
  % Bottom ports
  \foreach \pname/\pxoff/\pcol in {#8} {
    \pgfmathsetmacro{\px}{#2 + \pxoff}%
    \coordinate (#1-\pname) at (\px, \yB);
    \node[pinlabel, anchor=south, text=\pcol] at (\px, \yB + 0.10) {\detokenize\expandafter{\pname}};
  }
}
}

\begin{document}
\begin{tikzpicture}[
    >=Latex,
    font=\small,
    csrdatabox/.style={
        draw=hazardcolor, line width=1.5pt,
        fill=hazardcolor, fill opacity=0.15, text opacity=1.0,
        inner sep=0pt
    },
    stagetitle/.style={
        font=\bfseries\small, text=stagecolor, align=center
    },
    stageline/.style={dashed, draw=boundarycolor, line width=1.2pt},
    thickline/.style={draw=arrowgreen, line width=2.6pt, line cap=butt, line join=miter},
    thinline/.style={draw=arrowgreen, line width=1.2pt, line cap=butt, line join=miter},
    wirelabel/.style={font=\footnotesize\bfseries\ttfamily, text=arrowgreen}
]



% ECALL-only view: IF, ID and EX; MEM/WB do not participate in this path.
\foreach \xl/\xr/\title in {0/13/fetch\_stage,14.4/23.4/decode\_stage,24.8/36/execute\_stage} {
  \draw[draw=stagecolor,line width=1.8pt]
    (\xl,5.5) -- (\xr,5.5) -- (\xr,-3.3)
    decorate[decoration={zigzag,segment length=12pt,amplitude=3pt}] { -- (\xl,-3.3) } -- cycle;
  \node[anchor=north west,font=\bfseries\normalsize,text=stagecolor]
    at (\xl+0.3,5.35) {\texttt{\title}};
}
\foreach \x/\title in {6.5/Instruction Fetch (IF),18.9/Instruction Decode (ID),30.4/Execute (EX)}
  \node[stagetitle] at (\x,6) {\title};
\foreach \x in {13.35,23.75} {
  \draw[draw=regcolor,line width=1.3pt,fill=regcolor,fill opacity=0.15]
    (\x,5.5) -- (\x+0.7,5.5) -- (\x+0.7,-3.3)
    decorate[decoration={zigzag,segment length=8pt,amplitude=2.5pt}] { -- (\x,-3.3) } -- cycle;
}
\node[text=regcolor,font=\scriptsize] at (13.7,5.75) {IF/ID};
\node[text=regcolor,font=\scriptsize] at (24.1,5.75) {ID/EX};

% Blocks show only ports relevant to ECALL; PC selection is EX stage logic.
% Mux tapers toward its output; select enters above and MTVEC data enters from the left.
\draw[draw=stagecolor,line width=1.5pt,fill=stagecolor,fill opacity=0.15]
  (2,1.5) -- (8,0.8) -- (8,-0.8) -- (2,-1.5) -- cycle;
\node[font=\bfseries\normalsize,text=stagecolor] at (4.7,0.25) {next PC mux};
\coordinate (fetch-pc_src) at (5,1.15);
\coordinate (fetch-pc_next) at (8,0);
\coordinate (fetch-mtvec) at (2,-0.65);
\node[pinlabel,text=arrowgreen,anchor=north] at (5,1.05) {pc\_src};
\node[pinlabel,text=arrowgreen,anchor=east] at (7.88,0) {pc\_next};
\node[pinlabel,text=arrowpurple,anchor=west] at (2.12,-0.65) {MTVEC};
\drawblock{decode}{18.9}{0.55}{6.6}{Instruction\_Decode}%
  {instruction/arrowgreen}{is_ecall/arrowgreen}{}{2.1}
\drawblock{ecall}{30.1}{0}{7.6}{ecall\_processor}%
  {enable/arrowgreen, pc/arrowgreen}{}%
  {mepc_hw_request/-1.9/arrowpurple, mcause_hw_request/1.9/arrowpurple}{4}
\drawblock{select}{30.1}{3.8}{7.6}{PC source selection}%
  {is_ecall/arrowgreen}{pc_src/arrowgreen}{}{1.8}

% ECALL and its own PC are registered together at ID/EX.
\draw[thickline,->] (14.8,0.4) -- (decode-instruction);
\draw[thickline,->] (decode-is_ecall) -- (25.4,0.4) |- (ecall-enable);
\fill[arrowgreen] (25.4,0.4) circle (3pt);
\draw[thinline,->] (25.4,0.4) |- (select-is_ecall);
\draw[thickline,->] (21.8,-1.6) -- (25,-1.6) |- (ecall-pc);
\node[wirelabel,anchor=south east] at (23.3,-1.6) {pc\_cur};

% Combinational redirect back to IF, outside the forward pipeline registers.
\draw[thickline,->] (select-pc_src) -- (35,3.5) -- (35,7.2) -- (-0.5,7.2) -- (-0.5,3.6) -- (5,3.6) -- (fetch-pc_src);
\node[wirelabel,anchor=south] at (19.5,7.2) {ex\_to\_if.pc\_src};

\node[draw=regcolor,fill=regcolor,fill opacity=0.15,text=regcolor,text opacity=1,
  minimum width=2.2cm,minimum height=0.9cm] (pc-register) at (10.8,0) {PC register};
\draw[thickline,->] (fetch-pc_next) -- (pc-register.west);

% Top-level storage, in the same amber used for the hazard unit.
\draw[csrdatabox] (1,-6) rectangle (36,-10.7);
\node[anchor=north,font=\bfseries\normalsize,text=hazardcolor] at (20,-6.15) {csr\_data};
\foreach \xl/\xr/\name in {2/8/mtvec,24.8/30.4/mepc,31/35.4/mcause} {
  \draw[draw=hazardcolor,line width=1.5pt,fill=hazardcolor,fill opacity=0.15]
    (\xl,-7.2) rectangle (\xr,-9.6);
  \node[font=\bfseries,text=hazardcolor] at ({(\xl+\xr)/2},-8) {csr\_plugin};
  \node[font=\bfseries,text=hazardcolor] at ({(\xl+\xr)/2},-8.65) {\name};
}
\node[font=\footnotesize,text=arrowpurple] at (5,-9.2) {read\_data};
\node[font=\footnotesize,text=arrowpurple] at (27.6,-9.2) {write\_data = PC};
\node[font=\footnotesize,text=arrowpurple] at (33.2,-9.2) {write\_data = 11};
\draw[draw=arrowpurple,line width=2.6pt,->] (ecall-mepc_hw_request) -- (28.2,-7.2);
\draw[draw=arrowpurple,line width=2.6pt,->] (ecall-mcause_hw_request) -- (32,-7.2);
\node[font=\footnotesize\bfseries,text=arrowpurple,anchor=east] at (27.9,-4.7) {write\_enable = is\_ecall};

% Dedicated asynchronous read interface: dashed because RTL hookup is pending.
% No read-enable handshake exists in csr_hw_request_if.
\draw[draw=arrowpurple,line width=2.6pt,dashed,->]
  (5,-7.2) -- (5,-4.6) -- (0.7,-4.6) |- (fetch-mtvec);
\node[font=\footnotesize,text=arrowpurple,anchor=north,align=center] at (8,-11.1)
  {Dashed: intended mtvec hardware read connection (RTL hookup pending).\\
   Asynchronous read; no read-enable or request/response handshake.};
\node[font=\footnotesize,text=boundarycolor,anchor=north,align=center] at (28,-11.1)
  {ECALL path only; software CSR access and flush wiring omitted.\\
   Hardware writes take priority over software WB writes.};
\end{tikzpicture}
\end{document}
